\documentclass[12pt]{article}
\usepackage[T1]{fontenc}
\usepackage[utf8]{inputenc}
\usepackage[spanish]{babel}
\usepackage{graphicx}
\usepackage{subfigure}
\usepackage{amsmath}
\parindent = 0mm

\begin{document}
\title{Práctica	1. \\ Texto normal, párrafo,
Alineación y gráficos}
\author{Roly Sandro Gutierrez Benito}
% \date{} For no date
\maketitle
\tableofcontents

\section{Tipos y tamaños de fuente}
Iniciamos con la segunda práctica de este tutorial, 
colocando el título, autor y el día. Esta es la primera
sección donde conoceremos algunos tipos de letras, así
como los diferentes tamaños de letra que podemos usar.

\begin{itemize}
\item Tipo de letra \textbf{negrita}
\item Tipo de letra \textit{itálica}
\item Tipo de letra \textrm{romana}
\item Tipo de letra \textsf{sans serif}
\item Tipo de letra \textsl{inclinada}
\item Tipo de letra \textsc{Versalitas}
\end{itemize}

\newpage
\subsection{Tamaños de letra}
\begin{itemize}
\item {\tiny Tamaño} de letra
\item {\scriptsize Tamaño} de letra
\item {\footnotesize Tamaño} de letra
\item {\small Tamaño} de letra
\item {\normalsize Tamaño} de letra
\item {\large Tamaño} de letra
\item {\LARGE Tamaño} de letra
\item {\huge Tamaño} de letra
\item {\Huge Tamaño} de letra
\end{itemize}

\section{Párrafos, sangría y saltos de línea}
\subsection{Sangría}
Esta es la primera sección donde conoceremos algunos
tipos de letras, así como los diferentes tamaños de 
letra que podemos usar. Y ahora queremos también decidir
y modificar la sangría en el documento, aquí si iniciamos
con sagría.\\
\noindent Pero esta otra línea queremos que no tenga sangría.\\
Pero esta otra línea queremos que no tenga sangría.\\

\subsection{Salto de línea y nueva página}
Ahora también queremos decidir y modificar la sagría en el 
documento, aquí iniciamos con sangría.\\[2mm]
\par
Esta línea queremos que tenga un mayor salto de línea.

\subsection{Alineación de párrafos}
\begin{flushleft}
	Texto a la izquierda.
\end{flushleft}

\begin{center}
	Texto centrado
\end{center}

\begin{flushright}
	Texto a la derecha.
\end{flushright}

\subsection{Comillas y puntos suspensivos}
Este es un ejemplo: para comillas dobles,
`simples' y también ``dobles''. Y también los
puntos suspensivos \dots

\subsection{Espaciado horizontal}
Inicio \,fin \\[2mm]
Inicio \quad fin \\[2mm]
Inicio \qquad fin \\[2mm]
Inicio \hspace{3cm} fin \\[2mm]
Inicio \hspace{5cm} fin \\[2mm]
Inicio \hfill fin \\[2mm]

\subsection{Líneas de relleno}
Inicio \hrulefill \ fin\\[5mm]
Inicio \dotfill \ fin

\section{Incluir gráficas externas}
En \LaTeX podemos incluir figuras externas o 
generadas directamente en \LaTeX, utilizando algún
paquete, algunos de los formatos que podemos 
utilizar son \textbf{.jpg}, \textbf{.png},
\textbf{.eps} y \textbf{.pdf}.\\[5mm]

Los formatos vectoriales como lo son \textbf{.eps}
y \textbf{.pdf} son los más recomendados cuando
hacemos gráficas en las cuales queremos observar
detalles y presición, porque no pierden calidad al
aumentar o disminuir su tamaño, y para imágenes en
general podemos utilizar los formatos \textbf{.jpg}
o \textbf{.png}.

\section{Paquete graphicx}
\begin{center}
\includegraphics[scale=0.5]{figures/image1.eps}
\end{center}

\begin{itemize}
\item scale: Escala de la imagen
\item width: Ancho deseado de la imagen
\item height: Altura deseada de la imagen en cm.
\end{itemize}

\begin{center}
	\includegraphics[width=4cm]{figures/image1.eps}
\end{center}

\begin{center}
	\includegraphics[height=5cm]{figures/image1.eps}
\end{center}

\begin{center}
	\includegraphics[scale=0.1]{figures/image1.eps} \quad
	\includegraphics[scale=0.2]{figures/image1.eps} \quad
	\includegraphics[scale=0.3]{figures/image1.eps} \quad
	\includegraphics[scale=0.4]{figures/image1.eps} \quad
\end{center}

\newpage
\section{Objetos Flotantes}
Para poder ubicar una imagen en el lugar deseado se tiene
usar \textbf{!h} y dependiendo la posición que se quiere, se
puede aumentar: t (top), b (bottom)
\begin{figure}[!ht]
	\centering
	\includegraphics[scale=0.5]{figures/image1.eps}
	\caption{Imagen a escala 0.5}
	\label{figure1}
\end{figure}

\begin{itemize}
\item \textbf{t}: La imagen en la parte superior: (top)
\item \textbf{b}: La imagen en la parte inferior: (bottom)
\item \textbf{h}: La imagen en el lugar que se escribe: (here)
\end{itemize}

\begin{figure}[!ht]
	\centering
	\includegraphics[scale=1]{figures/image1.eps}
	\caption{Imagen eps}
	\label{figure2}
\end{figure}

\newpage
\subsection{Paquete subfigure}
\begin{figure}[!ht]
	\centering
	\subfigure[Primera figura]
	{\includegraphics[scale=0.05]{figures/image3.png}}
	\quad
	\subfigure[Segunda figura]
	{\includegraphics[scale=0.05]{figures/image3.png}}
	\quad
	\subfigure[Tercera figura]
	{\includegraphics[scale=0.05]{figures/image3.png}}
	\caption{Varias figuras}
	\label{figure3}
\end{figure}
Tal como aparece en la Figura\eqref{figure3}. Podemos
ver la forma en como poner tres imágenes en la misma fila.
\end{document}
